\documentclass[11pt,a4paper]{article}

\usepackage[utf8]{inputenc}
% Use Latin Modern where it is installed (Overleaf, full TeX Live, MiKTeX);
% otherwise fall back to the default Computer Modern fonts.
\IfFileExists{lmodern.sty}{\usepackage[T1]{fontenc}\usepackage{lmodern}}{}
\usepackage[margin=2cm]{geometry}
\usepackage{booktabs}
\usepackage{tabularx}
\usepackage{array}
\usepackage{amssymb}
\usepackage{enumitem}
\usepackage{xcolor}
\usepackage{listings}
\usepackage{titlesec}
\usepackage[colorlinks=true,linkcolor=accent,urlcolor=accent]{hyperref}

% ---------- Colours ----------
\definecolor{accent}{HTML}{1F4E79}
\definecolor{codebg}{HTML}{F5F7FA}
\definecolor{codekw}{HTML}{1F4E79}
\definecolor{rulegray}{HTML}{9AA5B1}

% ---------- Helpers ----------
\newcommand{\code}[1]{\texttt{#1}}
\newcolumntype{L}{>{\raggedright\arraybackslash}X}
\newcommand{\done}{$\square$}

\titleformat{\section}{\large\bfseries\color{accent}}{}{0pt}{}
\titlespacing*{\section}{0pt}{1.4em}{0.5em}

\setlength{\parindent}{0pt}
\setlength{\parskip}{0.5em}
\renewcommand{\arraystretch}{1.35}
\setlist[itemize]{leftmargin=1.4em, itemsep=0.2em, topsep=0.3em}

% ---------- SQL listings ----------
\lstset{
  language=SQL,
  basicstyle=\ttfamily\small,
  keywordstyle=\color{codekw}\bfseries,
  morekeywords={WITH, OVER, PARTITION, LAG, LEAD, ROW_NUMBER, RANK, DENSE_RANK, DATE_TRUNC, COALESCE, LIMIT},
  backgroundcolor=\color{codebg},
  frame=single,
  rulecolor=\color{rulegray},
  framesep=6pt,
  xleftmargin=8pt,
  xrightmargin=8pt,
  columns=fullflexible,
  keepspaces=true,
  showstringspaces=false
}

\begin{document}

% ---------- Title ----------
{\LARGE\bfseries\color{accent} SQL Roadmap for Data Science Interviews}\\[0.3em]
{\large 14 days, one hour a day, for someone who already thinks in pandas}\\[-0.4em]
{\color{rulegray}\rule{\textwidth}{0.6pt}}

Week 1 covers the basics. Week 2 covers what interviewers test most: CTEs,
window functions and dates, applied to fraud-style questions.

% ---------- Setup ----------
\section{Setup (Day 0, 20 minutes)}

\begin{itemize}
  \item \textbf{Dialect:} learn PostgreSQL syntax, which is the safest default for interviews.
  \item \textbf{Where to practise:} install DuckDB (\code{pip install duckdb}). It runs SQL
        directly on pandas DataFrames, so you can write a query and check it against the
        pandas version you already know.
  \item \textbf{Daily format:} 20 minutes on the concept, then 40 minutes typing queries yourself.
\end{itemize}

% ---------- Week 1 ----------
\section{Week 1: the basics}

{\small
\begin{tabularx}{\textwidth}{@{}c c L L L@{}}
\toprule
 & \textbf{Day} & \textbf{Topic} & \textbf{pandas equivalent} & \textbf{Watch out for} \\
\midrule
\done & 1 & \code{SELECT}, \code{WHERE}, \code{ORDER BY}, \code{LIMIT}, \code{DISTINCT}
      & column selection, boolean masks, \code{sort\_values}
      & SQL describes the result, not the steps \\
\done & 2 & \code{AND}/\code{OR}, \code{IN}, \code{BETWEEN}, \code{LIKE}, \code{IS NULL}, \code{COALESCE}
      & \code{isin}, \code{between}, \code{str.contains}, \code{fillna}
      & \code{col = NULL} is never true, so use \code{IS NULL} \\
\done & 3 & \code{COUNT}, \code{SUM}, \code{AVG}, \code{MIN}/\code{MAX}, \code{GROUP BY}
      & \code{groupby().agg()}
      & \code{COUNT(*)}, \code{COUNT(col)} and \code{COUNT(DISTINCT col)} differ \\
\done & 4 & \code{HAVING}, \code{CASE WHEN}, conditional aggregation
      & filtering after \code{groupby}, \code{np.where}
      & \code{WHERE} filters rows, \code{HAVING} filters groups \\
\done & 5 & \code{INNER JOIN}, \code{LEFT JOIN}
      & \code{merge(how=...)}
      & Check the row count before and after every join \\
\done & 6 & Self-joins, anti-joins, duplicate keys
      & \code{merge} on the same frame, \code{indicator=True}
      & Duplicate keys multiply rows, and NULL keys never match \\
\done & 7 & Review with 8 to 10 easy problems
      & ---
      & Execution order: \code{FROM}, \code{WHERE}, \code{GROUP BY}, \code{HAVING}, \code{SELECT}, \code{ORDER BY} \\
\bottomrule
\end{tabularx}\par
}

\textbf{Checkpoint for day 4.} You should be able to write fraud rate by merchant from memory:

\begin{lstlisting}
SELECT merchant_id,
       AVG(CASE WHEN is_fraud THEN 1.0 ELSE 0 END) AS fraud_rate
FROM transactions
GROUP BY merchant_id
HAVING COUNT(*) >= 100;
\end{lstlisting}

% ---------- Week 2 ----------
\section{Week 2: interview level}

{\small
\begin{tabularx}{\textwidth}{@{}c c L L L@{}}
\toprule
 & \textbf{Day} & \textbf{Topic} & \textbf{pandas equivalent} & \textbf{Practice question} \\
\midrule
\done & 8  & Subqueries and CTEs (\code{WITH})
      & intermediate DataFrames
      & Users whose spend is above the average \\
\done & 9  & \code{ROW\_NUMBER}, \code{RANK}, \code{DENSE\_RANK}
      & \code{groupby().cumcount()}, \code{rank}
      & Each user's first transaction \\
\done & 10 & \code{LAG}, \code{LEAD}
      & \code{groupby().shift()}
      & Gap since each user's previous transaction \\
\done & 11 & Running totals and window frames
      & \code{cumsum}, \code{rolling}
      & Running spend per user \\
\done & 12 & \code{DATE\_TRUNC}, \code{EXTRACT}, intervals
      & \code{dt.floor}, \code{dt.hour}, \code{Timedelta}
      & Fraud rate by week \\
\done & 13 & Fraud patterns combining all of the above
      & ---
      & Three or more transactions within ten minutes \\
\done & 14 & Timed mock: 5 medium problems, explained aloud
      & ---
      & --- \\
\bottomrule
\end{tabularx}\par
}

\textbf{Checkpoint for day 10.} Gap since each user's previous transaction:

\begin{lstlisting}
SELECT user_id,
       txn_time,
       txn_time - LAG(txn_time) OVER (
           PARTITION BY user_id
           ORDER BY txn_time
       ) AS gap_since_previous
FROM transactions;
\end{lstlisting}

% ---------- After ----------
\section{After day 14}

\begin{itemize}
  \item Do 30 to 50 problems across LeetCode's SQL set, DataLemur and StrataScratch,
        at 3 or 4 a day.
  \item Redo any problem you got wrong two days later, without looking at the solution.
  \item Prioritise fraud-style questions: transactions per user in the last 24 hours,
        first transaction per user, bursts of activity in a short window, and fraud rate
        by merchant or by week.
\end{itemize}

\end{document}